\documentclass[11pt]{article}
\usepackage[a4paper,margin=1in]{geometry}
\usepackage{amsmath,amssymb,amsthm,mathtools}
\usepackage{microtype}
\usepackage{graphicx}
\usepackage{hyperref}
\hypersetup{colorlinks=true,linkcolor=blue,citecolor=blue,urlcolor=blue}

\title{\textbf{Local Normal Forms for Complex Base Expansions:\\
A General $(s,b,l)$ Framework and the $\varphi$ Specialization}}
\author{VR \& AI}
\date{\today}

% --- theorem envs ---
\newtheorem{theorem}{Theorem}[section]
\newtheorem{lemma}[theorem]{Lemma}
\newtheorem{proposition}[theorem]{Proposition}
\theoremstyle{definition}
\newtheorem{definition}[theorem]{Definition}
\theoremstyle{remark}
\newtheorem{remark}[theorem]{Remark}
\newtheorem{example}[theorem]{Example}

% --- macros ---
\newcommand{\Val}{\mathrm{Val}}
\newcommand{\Z}{\mathbb{Z}}
\newcommand{\N}{\mathbb{N}}

\begin{document}
\maketitle

\section*{Notation crib (one page)}
\begin{tabular}{ll}
$l\ge2$ & sites per turn (angular radix)\\
$U=e^{i2\pi/l}$ & unit rotation per site\\
$b>1$ & radial growth per site\\
$B=b\,U$ & complex base\\
$\beta_l=b^l$ & full-turn radial factor on one lattice\\
$P,Q\in\mathbb Z$ & $b$ satisfies $b^2-Pb+Q=0$ (quadratic)\\
$U_n,V_n$ & Lucas sequences for $(P,Q)$, $X_{n+1}=P X_n-Q X_{n-1}$\\
$V_l=V_l(P,Q)$ & per-track digit-set size\\
$k=r+lt$ & track index ($r\in\{0,\dots,l-1\}$, $t\in\mathbb Z$)\\
$d_t$ & digit on a fixed track at site $t$\\
$\mathcal D$ & target digit set (nonnegative or balanced)\\
$E_q,\ E_q^{\mathrm{lin}}$ & weighted $\ell^1$ and signed linear potentials\\
\end{tabular}

\begin{abstract}
We present a local, phase-aware normalization theory for complex base expansions built from three parameters: sites per turn $l\in\N_{\ge2}$, per-step radial growth $b>1$, and a per-level shrinkage factor $s\in(0,1]$. For canonical integer-digit normalization we assume $b$ is a quadratic algebraic integer. We prove a \emph{tri-site, same-phase} carry identity that yields termination and confluence on each residue class (``phase track''), provide digit-set choices and carry-count bounds (with explicit denominator conditions), and show how to obtain global self-similarity with track-dependent angular lattices. We specialize to $\varphi$, recovering Lucas-number coefficients. We correct the termination proof to use the \emph{signed} $Q^{\,l}$ and provide an appendix covering the $Q^{\,l}<0$ case.
\end{abstract}

\section{Setup: the general $(s,b,l)$ model}\label{sec:setup}
Fix integers $l\ge2$ and parameters $b>1$, $s\in(0,1]$. Let
\[
U:=e^{i\,2\pi/l}\qquad\text{and}\qquad B:=b\,U.
\]

\begin{definition}[Complex $(b,l)$-string; phase tracks]\label{def:string}
A complex $(b,l)$-string is an integer sequence $(d_k)_{k\in\Z}$ satisfying:
\begin{enumerate}
\item \emph{finite right tail:} $\exists K\in\Z$ such that $d_k=0$ for all $k>K$;\vspace{2pt}
\item \emph{absolutely summable left tail:} $\sum_{t\ge0} |d_{-t}|\,b^{-t}<\infty$.
\end{enumerate}
Its value is the absolutely convergent series
\[
\Val(d):=\sum_{k\in\Z} d_k\,B^k,\qquad\text{with }|B|=b.
\]
Indices decompose into $l$ \emph{phase tracks} by residue $r\equiv k\pmod l$, where $k=r+lt$ and $t\in\Z$. On a fixed track $r$ we write $d_t:=d_{r+lt}$ and
\(
B^{r+lt}=U^r\,b^{r}\,\beta_l^t
\)
with $\beta_l=b^l>1$.
\end{definition}

\begin{definition}[Hierarchical sum with shrink $s$]\label{def:hier}
For levels $\ell=0,1,\dots$, a finite-depth address chooses integers $d_{\ell,k}$ and positions $k_\ell$; its value is
\[
z \;=\; \sum_{\ell\ge0} s^{\ell}\!\sum_{k\in\Z} d_{\ell,k}\,B^{k}.
\]
For infinite depth, $s^\ell\to0$ and absolute convergence of each inner sum ensure $z$ is well-defined.
\end{definition}

\begin{remark}[Quadratic $b$ for integer carries]\label{rem:quad}
To obtain \emph{integer} carry coefficients (and thus integer digits under normalization), we assume $b$ is a root of
\begin{equation}\label{eq:quad}
b^2 - P\,b + Q \;=\; 0,\qquad P,Q\in\Z,\ \ b>1.
\end{equation}
This includes all metallic means ($Q=-1$, $P\ge1$) and their integer powers $b^h$.
\end{remark}

\subsection*{Tracks and digit alphabets (notation)}
\emph{Phase tracks.} Indices decompose as $k=r+lt$ with $r\in\{0,\dots,l-1\}$ and $t\in\mathbb Z$.
On a fixed track $r$ we write $d_t:=d_{r+lt}$ and
\[
B^{r+lt} = U^r\,b^{r}\,\beta_l^{\,t},\qquad \beta_l:=b^l>1.
\]
\emph{Digit alphabets.} Throughout we normalize per track into a contiguous set $\mathcal D$ of size $V_l(P,Q)$:
\[
\mathcal D=
\begin{cases}
\{0,1,\dots,V_l-1\}, & \text{nonnegative},\\[2pt]
\{-a,\dots,a\},\quad a=\big\lfloor (V_l-1)/2\big\rfloor, & \text{balanced}.
\end{cases}
\]
Balanced digits are convenient when $Q^l<0$ (odd $l$), as they avoid long borrow cascades; nonnegative digits are natural when $Q^l>0$ (even $l$).

\subsection*{Convergence of hierarchical sums (sufficient bound)}
Let $\beta_l:=b^l>1$. If, for each level $\ell$, the per-track reindexed series obeys
\[
\sum_{t\in\Z}\! |d_{\ell,r+lt}|\,\beta_l^{t}\ \le\ C_\ell\quad\text{for all residues }r,
\]
then
\[
|z|\ \le\ \sum_{\ell\ge0} s^\ell \sum_{r=0}^{l-1} C_\ell.
\]
Thus $s<1$ and $\sup_\ell C_\ell<\infty$ imply absolute convergence.

\paragraph{Post-normalization quantification.}
After per-track normalization into a digit set of size $V:=V_l(P,Q)$, each digit satisfies $|d_{\ell,r+lt}|\le V-1$ (or $\le a=\lfloor(V-1)/2\rfloor$ in the balanced case). If the normalized support on level $\ell$ has $\nu_\ell$ nonzero sites in total and the maximal occupied site index is $t_{\max,\ell}$ (measured on its track), then the crude but effective bound
\[
\boxed{\ C_\ell\ \le\ (V-1)\,\beta_l^{\,t_{\max,\ell}}\,\nu_\ell\ }
\]
holds, since $\beta_l^{t}\le \beta_l^{t_{\max,\ell}}$ for all occupied $t$. When occupied sites form blocks, the sharper geometric sum
$\sum_{t=a}^{b}\beta_l^t=\beta_l^{a}\frac{\beta_l^{b-a+1}-1}{\beta_l-1}$
may be used trackwise; either way $C_\ell$ scales linearly in the (finite) active support at level $\ell$.

\paragraph{Numerical illustration ($\varphi,\,l=4$).}
Here $\beta_l=\varphi^4\approx 6.854$. Suppose a level $\ell$ has $\nu_\ell=3$ occupied sites on a track, at $t\in\{-1,0,1\}$, with digits bounded by $V-1=L_4-1=6$. The crude bound yields
\[
C_\ell \ \le\ (V-1)\,\beta_l^{\,t_{\max,\ell}}\,\nu_\ell \ =\ 6\cdot 6.854^{\,1}\cdot 3\ \approx\ 123.4.
\]
Using the geometric sum per block (here contiguous from $-1$ to $1$):
\[
\sum_{t=-1}^{1} \beta_l^{t} \ =\ \beta_l^{-1}\,\frac{\beta_l^{3}-1}{\beta_l-1}
\ \approx\ 0.146\cdot \frac{321.996-1}{5.854} \ \approx\ 8.0,
\]
so $C_\ell\le 6\cdot 8.0\approx 48$, a tighter bound.

\section{Local identity and normalization on one lattice $(b,l)$}\label{sec:one-lattice}

Let $U_n(P,Q),V_n(P,Q)$ be the Lucas sequences:
\[
U_0=0,\,U_1=1,\quad V_0=2,\,V_1=P,\quad X_{n+1}=P X_n - Q X_{n-1}.
\]
From \eqref{eq:quad} one has for $n\ge1$:
\begin{equation}\label{eq:lucas-power}
b^{2n}=V_n(P,Q)\,b^n - Q^n.
\end{equation}

\begin{proposition}[Phase-aligned (tri-site) carry identity]\label{prop:carry-general}
For all $k\equiv r\ (\mathrm{mod}\ l)$,
\begin{equation}\label{eq:carry-general}
B^{k+2l} - V_l(P,Q)\,B^{k+l} + Q^{l}\,B^{k} \;=\; 0.
\end{equation}
Equivalently on a track (with $\beta_l=b^l>1$):
\[
\beta_l^{t+1} \;=\; V_l(P,Q)\,\beta_l^t \;-\; Q^l\,\beta_l^{t-1}.
\]
\end{proposition}

\begin{proof}[Proof of Proposition~\ref{prop:carry-general}]
Let $b$ be quadratic with $b^2-Pb+Q=0$, and let $U=e^{i2\pi/l}$. The Lucas identity
$b^{2n}=V_n(P,Q)\,b^{n}-Q^{n}$ holds for all $n\ge1$. Taking $n=l$ gives
$b^{2l}=V_l(P,Q)\,b^{l}-Q^l$. Multiply by $U^{k}$ and note $U^{l}=1$:
\[
(bU)^{k+2l} - V_l(P,Q)\,(bU)^{k+l} + Q^{l}\,(bU)^k \;=\; 0.
\]
Since $B=bU$, this is exactly $B^{k+2l} - V_l(P,Q)\,B^{k+l} + Q^{l}B^k=0$, and it holds for every $k\equiv r\ (\mathrm{mod}\ l)$, i.e.\ along each residue track.
\end{proof}

\paragraph{Local moves (per track).}
For $t^\star\in\N$, define the value-preserving carry/borrow at center $t$ by
\begin{align}
\textbf{carry: }&(d_{t-1},d_t,d_{t+1})\mapsto (d_{t-1}+Q^{l}t^\star,\ d_t-V_l(P,Q)t^\star,\ d_{t+1}+t^\star),\label{eq:carry}\\
\textbf{borrow: }&(d_{t-1},d_t,d_{t+1})\mapsto (d_{t-1}-Q^{l}t^\star,\ d_t+V_l(P,Q)t^\star,\ d_{t+1}-t^\star).\nonumber
\end{align}

\subsection*{Algorithm 1: Deterministic per-track normalization}
Let $V:=V_l(P,Q)$ and let the target digit set be either $\{0,\dots,V-1\}$ (nonnegative) or the balanced set $\{-a,\dots,a\}$ with $a=\lfloor(V-1)/2\rfloor$. The following rightmost‑first procedure yields the canonical normal form.

\medskip
\noindent\textbf{Input:} integer array $(d_t)_{t\in\Z}$ with finite right tail. \quad
\textbf{Output:} normalized digits $(\hat d_t)$ in $\mathcal D$.

\begin{enumerate}
\item \textbf{Rightmost selection.} While some $t$ has $d_t\notin\mathcal D$, choose the \emph{maximum} such index $t^\star$.
\item \textbf{One-shot fix at $t^\star$.}
\begin{itemize}
  \item \emph{Nonnegative digits:} If $d_{t^\star}\ge V$, set $t=\big\lfloor d_{t^\star}/V\big\rfloor$ and apply a carry with multiplicity $t$:
  \[
  (d_{t^\star-1},d_{t^\star},d_{t^\star+1}) \leftarrow (d_{t^\star-1}+Q^l t,\ d_{t^\star}-Vt,\ d_{t^\star+1}+t).
  \]
  If $d_{t^\star}<0$, set $t=\big\lceil(-d_{t^\star})/V\big\rceil$ and apply a borrow of multiplicity $t$:
  \[
  (d_{t^\star-1},d_{t^\star},d_{t^\star+1}) \leftarrow (d_{t^\star-1}-Q^l t,\ d_{t^\star}+Vt,\ d_{t^\star+1}-t).
  \]
  \item \emph{Balanced digits:} Replace the division by $V$ with centered rounding $t=\big\lfloor (d_{t^\star}+a)/V\big\rfloor$ so that the new center lands in $\{-a,\dots,a\}$ in one step.
\end{itemize}
\item \textbf{Repeat} Step~1.
\end{enumerate}

\paragraph{Remark (Sparse / NAF-like variants).}
If a sparser normal form is desired (e.g., forbidding adjacent nonzeros on a track, or minimizing a local adjacency cost), two light-touch adaptations keep the theory intact:
(i) a rightmost-first selection with a tie-breaker that reduces the adjacency count locally; and
(ii) a bounded post-normalization sparsify pass (balanced digits) that removes adjacent pairs via a single local carry/borrow while keeping digits in $\mathcal D$.
Both preserve termination and confluence.

\begin{theorem}[Termination and confluence (single lattice)]\label{thm:terminate}
Fix a track and $q\in(1,\beta_l)$ with $\beta_l=b^l$. With the \emph{signed} linear potential $E_q^{\mathrm{lin}}=\sum_t d_t q^t$, a unit carry at $t$ changes $E_q^{\mathrm{lin}}$ by
$q^{t-1}(q^2 - V_l q + Q^l)$, which is negative for all $q\in(1,\beta_l)$.
Hence normalization terminates on the track. Overlaps of carries occur only at adjacent centers; critical pairs are joinable. Summing the $l$ tracks yields a unique normal form.
\end{theorem}

\begin{proof}[Termination]
A unit carry at $t$ changes $(d_{t-1},d_t,d_{t+1})$ by $(+Q^l,\,-V_l,\,+1)$, hence
\(
\Delta E_q^{\mathrm{lin}}=q^{t-1}(q^2 - V_l q + Q^l)=(q-\beta_l)(q-Q^l/\beta_l)\,q^{t-1}<0
\)
for $q\in(1,\beta_l)$ since $Q^l/\beta_l<1<\beta_l$.
For the nonnegative target with $Q^l>0$ (e.g.\ metallic means and even $l$), all steps are carries so $E_q^{\mathrm{lin}}\ge0$ strictly decreases.
For $Q^l<0$ see Appendix~A.
\end{proof}

\begin{lemma}[Critical-pair check (adjacent carries commute)]\label{lem:commute}
Let $C_t$ denote a unit carry at center $t$ (i.e.\ $t^\star=1$), and write $V:=V_l(P,Q)$, $Q_0:=Q^l$.
Restrict to $\{t-1,t,t+1,t+2\}$ with digits $(a,b,c,d)$. Then
\[
C_{t+1}\,C_t(a,b,c,d)\ =\ C_t\,C_{t+1}(a,b,c,d)\ =\ (a+Q_0,\ b-V+Q_0,\ c+1-V,\ d+1).
\]
By linearity, carries with $t^\star>1$ decompose into sums of unit carries. Hence overlapping carries at $t$ and $t{+}1$ commute.
\end{lemma}

\begin{proof}
Apply the updates in both orders and compare; the two linear maps coincide on $(a,b,c,d)$.
\end{proof}

\begin{proof}[Confluence]
Local confluence follows from commutation of adjacent carries and disjointness of nonadjacent ones. With termination, Newman's lemma yields global confluence and a unique per-track normal form.
\end{proof}

\section{Digit-set variants and fast addition}\label{sec:digits-fast}
\begin{itemize}
\item \textbf{Balanced signed digits.} With $\mathcal D=\{-a,\dots,a\}$, a provisional sum $s_t$ satisfies $|s_t|\le 2a\le V_l-1$. A \emph{single} carry at $t$ reduces $|s_t|$ to $\le a$; this yields one-pass, bounded-carry addition per track (plus bounded neighbor fixes).
\item \textbf{Sparse (NAF-like).} Use a small signed set (e.g.\ $\{-1,0,+1\}$) and forbid adjacent nonzeros on a track. Add a local penalty to the potential to keep termination; sparsity reduces Hamming weight.
\end{itemize}

\section{Carry-count bounds (per track)}\label{sec:bounds}
All bounds below presuppose the denominators are positive. In particular, the uniform $\ell^1$ estimate requires $V_l(P,Q)>(|Q|^{l}+1)$; when this gap is small, prefer the weighted bound.

\begin{proposition}[Uniform $\ell^1$ bound]
Assume $V_l(P,Q)>(|Q|^{l}+1)$. Each carry reduces the center by $V_l(P,Q)$ and increases the two neighbors' absolute sum by at most $|Q|^{l}+1$. Hence
\[
\#\mathrm{carries}\ \le\ \frac{S}{\,V_l(P,Q) - (|Q|^{l}+1)\,}.
\]
\end{proposition}

\begin{proposition}[Weighted bound]
For $q\in(1,b^l)$ and $E_q=\sum_t |d_t|q^t$, if $V_l(P,Q)^2/4 - |Q|^{l}>0$ (e.g.\ $q=V_l/2<b^l$), then
\[
\#\mathrm{carries}\ \le\ \frac{E_q^{\mathrm{init}}}{\,\tfrac{V_l(P,Q)^2}{4}-|Q|^{l}\,}.
\]
\end{proposition}

\paragraph{Carry-count instantiation ($\varphi,\,l=4$).}
On one track with initial digits $(0,13,12,0)$ we have $S=25$, $V=7$, $|Q|^l=1$.
The uniform $\ell^1$ bound gives $\#\le 25/(7-2)=5$.
With $q=3.5$, $E_q=13+12q=55$, so the weighted bound gives $\#\le 55/(49/4-1)\approx 4.89$.
A concrete normalization needs exactly 4 carries, showing both bounds are nontrivial and reasonably tight.

\subsection*{Bound tightness across configurations}
We compare the uniform $\ell^1$ bound and the weighted bound on several initial states for $(\varphi,l{=}4)$, using rightmost-first normalization with unit-carry counting.\footnote{Actual carry counts were computed by Algorithm~1, counting a carry of multiplicity $t$ as $t$ unit tri-site carries.}

\begin{center}
\begin{tabular}{lrrrrr}
\hline
case & $S$ & $E_q$ ($q{=}3.5$) & uniform & weighted & actual \# \\
\hline
A\; sparse center $(0,13,0)$ & 13 & 13.000 & 2.600 & 1.156 & 1 \\
B\; adjacent heavy $(0,13,12,0)$ & 25 & 55.000 & 5.000 & 4.889 & 4 \\
C\; dense block $(12,12,12,12,12)$ & 60 & 205.408 & 12.000 & 18.259 & 10 \\
D\; far-right spike $(0,0,0,0,50)$ & 50 & 612.500 & 10.000 & 54.444 & 9 \\
\hline
\end{tabular}
\end{center}

\section{Track-dependent $l_r$ with global self-similarity}\label{sec:selfsimilar}
Fix a \emph{global similarity exponent} $M\in\N$ and set $\beta:=b^{M}>1$. Choose tracks $r$; for each track pick $l_r\mid M$ and define
\[
h_r:=\frac{M}{l_r}\in\N,\qquad B_r:=b^{h_r}\,e^{i2\pi/l_r}.
\]
Then $(b^{h_r})^{l_r}=\beta$ (same full-turn radial factor on all tracks). With level shrink $s^\ell$, the value is
\[
z=\sum_{\ell\ge0}s^\ell \sum_{r}\sum_{k\in\Z} d_{r,\ell,k}\,B_r^{k}.
\]

\begin{proposition}[Uniform carry across tracks]\label{prop:uniform}
If $b$ satisfies $x^2-Px+Q=0$, then for each track $r$ and $k\equiv r\ (\mathrm{mod}\ l_r)$,
\begin{equation}\label{eq:uniform}
B_r^{k+2l_r}-V_{M}(P,Q)\,B_r^{k+l_r}+Q^{M}\,B_r^{k}=0.
\end{equation}
Thus every track uses the \emph{same} coefficients $V_{M}(P,Q)$ and $Q^{M}$; the digit-set size is $V_{M}(P,Q)$. Termination/confluence and bounds hold uniformly with $\beta=b^{M}$.
\end{proposition}

\begin{proof}[Proof of Proposition~\ref{prop:uniform}]
Fix the global $M\in\N$ and a track with $l_r\mid M$, $h_r=M/l_r$. Then $b^{h_r}$ is quadratic over $\Z$ with the same $(P,Q)$, hence
\[
(b^{h_r})^{2l_r} = V_{l_r}(P,Q)\,(b^{h_r})^{l_r} - Q^{\,l_r}.
\]
Because $h_r l_r=M$, the composed Lucas coefficient is $V_{h_r l_r}(P,Q)=V_M(P,Q)$ and $Q^{h_r l_r}=Q^{M}$. Thus
$b^{2M}=V_M(P,Q)\,b^{M}-Q^{M}$. Multiplying by $U_r^{k}$ with $U_r^{l_r}=1$ yields
\[
B_r^{k+2l_r} - V_M(P,Q)\,B_r^{k+l_r} + Q^{M}B_r^{k}=0,
\]
which is track‑independent.
\end{proof}

\section{Specialization to $\boldsymbol\varphi$}\label{sec:phi}
Set $b=\varphi=\tfrac{1+\sqrt5}{2}$, which satisfies $x^2-x-1=0$ ($P=1$, $Q=-1$). Then
\[
V_n(P,Q)=L_n\quad\text{(Lucas numbers)},\qquad Q^{n}=(-1)^n.
\]
\paragraph{Single lattice $(\varphi,l)$.}
The carry identity becomes
\[
B^{k+2l}-L_l\,B^{k+l}+(-1)^l B^k=0,\qquad B=\varphi\,e^{i2\pi/l},
\]
and per-track digits normalize into a set of size $L_l$; with $l$ even, $\{0,\dots,L_l-1\}$ is natural. The potential window is $(1,\varphi^{l})$.

\begin{example}[$\varphi$ and quarter-turns]
Take $l=4$. Then $L_4=7$ and $B^{k+8}-7B^{k+4}+B^k=0$ ($(-1)^4=+1$).
Nonzeros, e.g.\ at $k\in\{-2,2,5\}$, decompose by residues:
$r=2$ has $t=-1$ and $t=0$; $r=1$ has $t=1$; tracks $r=0,3$ empty.
Since digits lie in $\{0,\dots,6\}$, no carries are needed here.
\end{example}

\section{Figures (selected)}
% Figures referenced earlier in the project; keep includes if desired
% \includegraphics{fig_spiral_rays_phi_l4.png} etc.

\section{Related work (brief pointers)}
Complex bases and non-standard radices have a rich history.
Knuth \cite{Knuth1997} discusses negative and unusual bases in a foundational algorithmic context.
Penney \cite{Penney1965} introduced base $-1+i$ arithmetic.
Gilbert \cite{Gilbert1981,Gilbert1983} studied arithmetic and representations in complex bases.
Frougny and Sakarovitch \cite{FrougnySakarovitch1992} surveyed automatic/normalized representations in algebraic bases.
Our contribution is a \emph{phase-aware tri-site} normalization that (i) works uniformly on residue tracks, (ii) gives simple termination/confluence proofs via a potential, and (iii) extends cleanly to track-dependent lattices with \emph{uniform} carry coefficients under a global self-similarity constraint
(see also \cite{Heuberger2007,AkiyamaPetho2011} for modern developments on canonical number systems in quadratic/imaginary settings).

\section*{Appendix A: Termination for $Q^{\,l}<0$ (sketch)}
Two variants keep the Lyapunov argument valid for odd $l$ (e.g.\ metallic means).
\begin{enumerate}
\item \textbf{Balanced digits (macro-steps).}
Normalize into $\mathcal D=\{-a,\dots,a\}$. Group updates into macro-steps:
perform a rightmost carry at $t$ (unit multiplicity) and, if $d_{t-1}$ leaves $\mathcal D$, immediately fix it by at most one local borrow.
The signed linear potential $E_q^{\mathrm{lin}}$ decreases by $q^{t-1}(q^2-V_l q + Q^l)$ per macro-step for any $q\in(1,\beta_l)$.
\item \textbf{Two-sided Lyapunov.}
Use $E_q^\lambda=E_q^{\mathrm{lin}}+\lambda\sum_t |d_t|q^t$ with small $\lambda>0$.
Carries decrease $E_q^{\mathrm{lin}}$ by a fixed amount; transient borrows change the $\ell^1$ term by $O(q^{t-1})$ with a controlled coefficient. Choose $\lambda$ so the net change is negative.
\end{enumerate}
Local confluence remains identical (adjacent tri-site updates commute), so Newman’s lemma yields uniqueness.

\section*{Appendix B: Tiny reference implementation (for $\varphi,l=4$)}
A 30-line Python script (rightmost-first, nonnegative digits $\{0,\dots,6\}$) for per-track normalization and bound checks is provided as supplementary material (\texttt{phi\_l4\_reference.py}). Core loop:
\begin{verbatim}
V = 7; Qpow = 1
def normalize_rightmost_first(d):
    d = d.copy(); unit_carries = 0
    while True:
        cand = [t for t,v in d.items() if v >= V]
        if not cand: break
        tstar = max(cand); tmult = d[tstar] // V
        d[tstar-1] = d.get(tstar-1, 0) + Qpow * tmult
        d[tstar]   = d[tstar] - V * tmult
        d[tstar+1] = d.get(tstar+1, 0) + tmult
        unit_carries += tmult
    return unit_carries, d
\end{verbatim}

\bibliographystyle{plain}
\bibliography{refs}

\end{document}
